% Generated draft: requires mathematical and source review.
\documentclass[11pt]{article}
\usepackage[margin=25mm]{geometry}
\usepackage{fontspec}
\setmainfont{Latin Modern Roman}
\usepackage{amsmath,amssymb,mathtools,unicode-math}
\setmathfont{Latin Modern Math}
\usepackage{enumitem,needspace,xurl,hyperref}
\setlength{\parindent}{0pt}
\setlength{\parskip}{5pt}
\setlength{\emergencystretch}{3em}
\newcommand{\N}{\mathbb{N}}
\newcommand{\Z}{\mathbb{Z}}
\newcommand{\Q}{\mathbb{Q}}
\newcommand{\R}{\mathbb{R}}
\newcommand{\C}{\mathbb{C}}
\newcommand{\F}{\mathbb{F}}
\newcommand{\A}{\mathbb{A}}
\newcommand{\PP}{\mathbb{P}}
\newcommand{\E}{\mathbb{E}}
\newcommand{\eps}{\varepsilon}
\newcommand{\dd}{\,\mathrm{d}}
\begin{document}
\section*{UnsolvedMath ID: 30003278}
\textbf{Source problem code:} catalogue entry
\subsection{Definitions and mathematical statement}
The supplied record contains two separate questions.

\begin{enumerate}
\item With the conventional interpretation that \(p\) denotes a prime, let \(p>3\), let \(\F_p\) be the field with \(p\) elements, and write \(\F_p^\times=\F_p\setminus\{0\}\). For each \(t\in\F_p^\times\), define
\[
E_p(t)=\bigl\{(x,y)\in(\F_p^\times)^2:x+x^{-1}+y+y^{-1}=t\bigr\}.
\]
All operations in this definition take place in \(\F_p\). Does there exist, for every such \(p\) and \(t\), an explicit bijection
\[
\Phi_{p,t}:E_p(t)\longrightarrow E_p(16t^{-1})?
\]
Here \(16t^{-1}\) is a nonzero field element because \(p>3\) and \(t\ne0\). The source asks for an explicit bijection or a bijective proof, but does not specify what form of construction qualifies as explicit. Its background describes a point-count identity without supplying its proof or a construction of the requested bijection.

\item Define
\[
S=\bigl\{n\in\Z_{>0}:\text{every prime }q\text{ dividing }n\text{ satisfies }q\equiv3\pmod4\bigr\}.
\]
Thus \(1\in S\), since it has no prime divisors, and \(0\notin S\). Is the following assertion true?
\[
\exists N_0\in\Z_{>0}\ \forall N\in\Z_{>0},\qquad
\bigl(N\ge N_0\ \land\ N\equiv0\pmod2\bigr)
\Longrightarrow
\exists a,b\in S\quad N=a+b.
\]
The two summands are not required to be distinct.
\end{enumerate}

The supplied record labels the composite problem open. Its dated assessment specifically describes the second question as unresolved; this is supplied curation, not an independently verified conclusion about either question.
\subsection{Short English statement}
For every prime \(p>3\) and nonzero \(t\in\F_p\), construct an explicit bijection between the two specified sets of finite-field points, or determine whether such a construction exists. Separately, determine whether every sufficiently large positive even integer is a sum of two positive integers whose prime divisors are all congruent to \(3\) modulo \(4\). The required meaning of ``explicit'' is unspecified.
\subsection{Sources}
\begin{itemize}
\item Local UnsolvedMath registry, supplied complete source record: permanent ID \(30003278\), source problem code \texttt{catalogue entry}, in the Oberwolfach Reports Open Problems collection. The supplied proposer field is empty; no individual proposer is identified. The recorded proposed year is 2016.
\item Supplied report citation: \emph{Analytic Number Theory} (2017), \emph{Oberwolfach Reports}, pp.\ 2975--3030, DOI \texttt{10.4171/owr/2016/53}. Supplied source link: \url{https://doi.org/10.4171/owr/2016/53}.
\item Supplied literature source, labelled ``OWR 53 report'': \url{https://publications.mfo.de/bitstream/handle/mfo/3557/OWR_2016_53.pdf?isAllowed=y&sequence=1}.
\item Literature assessment dated exactly \texttt{2026-08-21}: the supplied curation labels the composite record open, describes a finite-field point-count identity, and says that the sufficiently-large-even-integer assertion remains unresolved. This is supplied dated curation and has not been independently verified here; time-sensitive status and cited claims require independent checking.
\item The supplied statement status is \texttt{corrected\_verified}. Its verification note reports that the official problem-session report was checked and that this record retains two source questions omitted from the overlapping composite extraction \texttt{OWR-15174-010}. These are supplied provenance claims, not verification performed for this draft.
\end{itemize}
\end{document}
